% =====================================================================
\section{The window in the van Trees bound}
\label{app:prior}
% =====================================================================

For a prior density $p(\theta)$, the van Trees inequality \cite{vantrees1968,gill1995} bounds the mean squared error
of any estimator, averaged over the prior, by
\begin{equation*}
  \sigma_\theta^2(n)\ \ge\ \big(J(p)+nI^*\big)^{-1},\qquad
  J(p)=\int\frac{p'(\theta)^2}{p(\theta)}\,d\theta .
\end{equation*}
For a Gaussian prior $J(p)=\sigma_0^{-2}$, the value used for tracking. A uniform window drops from $1/W$ to zero at
its edges, so $J(p)$ is infinite and the inequality reduces to $\sigma_\theta\ge0$: true, but empty. The inverse
variance $12/W^2$ cannot simply take its place, since $J(p)\ge1/\mathrm{Var}(\theta)$ with equality only for a
Gaussian: for a signal linear in $\theta$, $(12/W^2+nI^*)^{-1/2}$ is the error of the best linear estimator, which an
estimator can reach and the best estimator beats, so it is not a lower bound.

\paragraph{Worst-placed resonance.} Instead of the prior, weight the positions in the window by
$q(v)=(2/W)\cos^2(\pi v/W)$, with $v$ the distance from the center. Of all densities confined to the window this one
has the least Fisher information, $J(q)=4\pi^2/W^2$ \cite{gill1995}. The van Trees inequality then bounds the mean
squared error averaged with the weight $q$, and the worst-placed resonance is at least as bad as any such average,
which gives Eq.~\eqref{eq:vantrees} for a global search. The prior itself stays uniform; $q$ is only a step in the
argument.

\paragraph{Average over the window.} This argument does not cover the error averaged over a uniformly placed
resonance, which is what the runs report. For a signal linear in $\theta$ the best achievable average error can be
computed numerically. It lies above the bound until $nI^*W^2\approx400$, by a factor that starts at
$\pi/\sqrt3\approx1.8$ at $n=0$, and below it from there on, by at most $1.1\%$ in $\sigma$, an effect that vanishes as $n$ grows. The cause
is the hard edge: the posterior of a resonance close to an edge is cut off there, so it is slightly easier to locate.
